\documentclass[12pt]{article}

% =========================
% 0) Metadata (EDIT ME)
% =========================
\newcommand{\CourseCode}{MATH 431}
\newcommand{\CourseTitle}{Introduction to Probability}
\newcommand{\CourseSection}{LEC 002} % change to 001 or 003 if needed
\newcommand{\Semester}{Spring 2026}
\newcommand{\Instructor}{Sarah Tammen} % LEC 001: Ander Aguirre | LEC 002: Sarah Tammen | LEC 003: Nicholas Derr
\newcommand{\MeetingInfo}{MWF 9:55--10:45, 6102 Sewell Social Sciences} % LEC 001 default; update if in 002/003
\newcommand{\HomeworkNumber}{7}
\newcommand{\YourName}{Alejandro De La Torre}
\newcommand{\DueDate}{March 13, 2026} % or e.g. February 2, 2026

% =========================
% 1) Core math tools
% =========================
\usepackage{amsmath, amssymb, amsthm}

% =========================
% 2) Lists and spacing
% =========================
\usepackage{enumitem}
\setlist[enumerate,1]{label=\textbf{(\alph*)}, leftmargin=2em, itemsep=0.25em}
\setlist[itemize]{leftmargin=1.5em, itemsep=0.25em}
\setlength{\parskip}{0.35em}
\setlength{\parindent}{0pt}

% =========================
% 3) Page geometry + header/footer
% =========================
\usepackage{geometry}
\geometry{margin=1in}

\usepackage{fancyhdr}
\pagestyle{fancy}
\fancyhf{}
\lhead{\CourseCode\ --- \CourseSection, \Semester}
\rhead{Homework \HomeworkNumber}
\cfoot{\thepage}
\renewcommand{\headrulewidth}{0.4pt}
\setlength{\headheight}{15pt}
\addtolength{\topmargin}{-2pt}

% =========================
% 4) Links (subtle)
% =========================
\usepackage[hidelinks]{hyperref}
\setcounter{tocdepth}{1}

% =========================
% 5) Figures (optional)
% =========================
\usepackage{graphicx}
\graphicspath{{images/}}

% =========================
% 6) Lightweight boxes (optional)
% =========================
\usepackage{mdframed}
\newmdenv[
  skipabove=0.6\baselineskip,
  skipbelow=0.6\baselineskip,
  linewidth=0.6pt,
  roundcorner=4pt,
  innerleftmargin=10pt,
  innerrightmargin=10pt,
  innertopmargin=8pt,
  innerbottommargin=8pt
]{boxnote}

% =========================
% 7) Probability / notation macros
% =========================
\newcommand{\R}{\mathbb{R}}
\newcommand{\N}{\mathbb{N}}
\newcommand{\Z}{\mathbb{Z}}

\newcommand{\OmegaSpace}{\Omega}
\newcommand{\FField}{\mathcal{F}}
% Probability measure in case it is relevant or want to distinguish it from P(A)
%\newcommand{\PMeasure}{\mathbb{P}} 
\newcommand{\E}{\mathbb{E}}

\newcommand{\given}{\,\middle|\,}
\newcommand{\ind}{\mathbf{1}}

% Common “defined as” symbol
\newcommand{\defeq}{\vcentcolon=}

% =========================
% 8) Problem command (TOC + PDF bookmarks)
% Usage:
%   \problem{<label>}
%   \problem[Short title]{<label>}
% Example:
%   \problem{1.4}
%   \problem[Sample space]{1.4}
% =========================
\makeatletter
\newcommand{\problem}[2][]{%
  \def\prob@title{Problem #2\if\relax\detokenize{#1}\relax\else\ (\textit{#1})\fi}%
  \phantomsection%
  \pdfbookmark[1]{\prob@title}{prob#2}%
  \addcontentsline{toc}{section}{\prob@title}%
  \section*{\prob@title}%
  \vspace{-0.25em}\hrule\vspace{0.8em}%
}
\makeatother

% =========================
% 9) Solution environment
% =========================
\newenvironment{solution}{%
  \noindent\textbf{Solution.}\ \ignorespaces
}{\par}

% Optional scratch / planning block
\newenvironment{scratch}{%
  \begin{boxnote}\textbf{Scratch / Setup.}\ \small
}{\end{boxnote}}

% Final answer command: expects math only (no $$ or \[ \])
\newcommand{\finalanswer}[1]{%
  \par\medskip
  \noindent\textbf{Final:}\ \fbox{$#1$}\par
} % AGAIN: MATH ONLY INSIDE THE BOX; if word answer, use \text{} inside or use \fbox{TEXT} or \framebox

% =========================
% 10) Title block
% =========================
\title{Homework \HomeworkNumber}
\author{\YourName \\
\CourseCode\ --- \CourseTitle \\
\CourseSection \\
Instructor: \Instructor}
\date{Due: \DueDate}

\begin{document}
\maketitle

% \section*{Collaboration / Resources}
% (Optional) Write “I worked alone” or list collaborators/resources.

\tableofcontents
\bigskip
\hrule
\bigskip
\newpage

% ============================================================
% Problems (duplicate blocks as needed)
% ============================================================


\problem{4.4}
Hint. You will need the continuity correction for this.
Liz is standing on the real number line at position 0. She rolls a die repeatedly. If the roll is 1 or 2, she takes \emph{one} step to the right (in the positive direction). If the roll is 3, 4, 5 or 6, she takes \emph{two} steps to the right. Let $X_n$ be Liz's position after $n$ rolls of the die. Estimate the probability that $X_{90}$ is at least 160.

\begin{solution}
\begin{scratch}
Let $Y_i$ be the step on roll $i$: $Y_i=1$ with probability $1/3$ (roll 1 or 2) and $Y_i=2$ with probability $2/3$ (roll 3--6). Then $X_n=\sum_{i=1}^n Y_i$.
Let $S_n$ be the number of two-steps. Then $S_n\sim\text{Bin}(n,2/3)$ and $X_n = n + S_n$.
For $n=90$, $E[S_{90}]=60$, $\mathrm{Var}(S_{90})=20$, $\sigma=\sqrt{20}$.
We need $P(X_{90}\ge160)=P(S_{90}\ge70)$ and use normal approximation with continuity correction.
\end{scratch}

Since $X_{90}=90+S_{90}$,
\[
P(X_{90}\ge 160)=P(S_{90}\ge 70)\approx P\left(\frac{S_{90}-60}{\sqrt{20}}\ge \frac{69.5-60}{\sqrt{20}}\right).
\]
Thus
\[
P(X_{90}\ge 160)\approx 1-\Phi(2.12)\approx 0.0168.
\]

\finalanswer{P(X_{90}\ge160)\approx 0.0168}
\end{solution}

\newpage

\problem{4.6}
A pollster would like to estimate the fraction $p$ of people in a population who intend to vote for a particular candidate. How large must a random sample be in order to be at least 95\% certain that the fraction $\hat{p}$ of positive answers in the sample is within 0.02 of the true $p$?

\begin{solution}
\begin{scratch}
Let $S_n\sim\text{Bin}(n,p)$ be the number of positive answers and $\hat p=S_n/n$.
We want $P(|\hat p-p|\le 0.02)\ge 0.95$.
Use the normal approximation for $\hat p$: 
$P(|\hat p-p|<\varepsilon)\approx 2\Phi\!
\left(\dfrac{\varepsilon\sqrt{n}}{\sqrt{p(1-p)}}\right)-1$.
Since $p(1-p)\le 1/4$, a conservative bound is
$P(|\hat p-p|<\varepsilon)\gtrsim 2\Phi(2\varepsilon\sqrt{n})-1$.
\end{scratch}

With $\varepsilon=0.02$, require
\[
2\Phi(0.04\sqrt{n})-1\ge 0.95 \quad\Longrightarrow\quad \Phi(0.04\sqrt{n})\ge 0.975.
\]
Using $z_{0.975}=1.96$,
\[
0.04\sqrt{n}\ge 1.96 \quad\Longrightarrow\quad \sqrt{n}\ge 49 \quad\Longrightarrow\quad n\ge 2401.
\]

\finalanswer{n\ge 2401}
\end{solution}

\newpage

\problem{4.8}
In a million rolls of a biased die the number 6 shows up 180,000 times. Find a 99.9\% confidence interval for the unknown probability that the die rolls 6.

\begin{solution}
\begin{scratch}
Let $S\sim\text{Bin}(n,p)$ be the number of 6's in $n=1{,}000{,}000$ rolls and $\hat p=S/n=0.18$.
For a 99.9\% confidence interval, we want $P(|\hat p-p|<\varepsilon)\approx 0.999$.
Normal approximation gives
$P(|\hat p-p|<\varepsilon)\approx 2\Phi\!\left(\dfrac{\varepsilon\sqrt{n}}{\sqrt{p(1-p)}}\right)-1$.
Using $\sqrt{p(1-p)}\le 1/2$ yields the conservative bound
$P(|\hat p-p|<\varepsilon)\gtrsim 2\Phi(2\varepsilon\sqrt{n})-1$.
Set $2\Phi(2\varepsilon\sqrt{n})-1=0.999$ so $2\varepsilon\sqrt{n}=z_{0.9995}\approx 3.29$.
\end{scratch}

Thus
\[
\varepsilon\approx \frac{3.29}{2\sqrt{1{,}000{,}000}}=\frac{3.29}{2000}\approx 0.00165,
\]
and the 99.9\% confidence interval is
\[
\hat p\pm \varepsilon = 0.18\pm 0.00165=(0.17835,\,0.18165).
\]

\finalanswer{(0.17835,\,0.18165)}
\end{solution}

\newpage

\problem{4.10}
A hockey player scores at least one goal in roughly half of his games. How would you estimate the percentage of games where he scores a hat-trick (three goals)?

\begin{solution}
\begin{scratch}
Model goals per game by $G\sim\text{Poisson}(\lambda)$ (rare scoring chances). Then
$P(G\ge 1)=1-e^{-\lambda}\approx 0.5$ so $\lambda=\ln 2$.
A hat-trick means exactly three goals, so compute $P(G=3)$.
\end{scratch}

With $\lambda=\ln 2$,
\[
P(G=3)=e^{-\lambda}\frac{\lambda^3}{3!}=\frac{(\ln 2)^3}{6}e^{-\ln 2}
=\frac{(\ln 2)^3}{12}\approx 0.0278.
\]
So the estimate is about $2.8\%$ of games.

\finalanswer{P(\text{hat-trick})\approx 0.0278}
\end{solution}

\newpage

\problem{4.16}
We choose 500 numbers uniformly at random from the interval $[1.5, 4.8]$.
\begin{enumerate}[label=(\alph*)]
\item Approximate the probability of the event that less than 65 of the numbers start with the digit 1.
\item Approximate the probability of the event that more than 160 of the numbers start with the digit 3.
\end{enumerate}

\begin{solution}
\begin{scratch}
A number in $[1.5,4.8]$ starts with digit 1 iff it lies in $[1.5,2)$; length $0.5$.
So $p_1=0.5/(4.8-1.5)=5/33$. Let $S\sim\text{Bin}(500,p_1)$ be the count.
For digit 3, the interval is $[3,4)$ of length $1$, so $p_3=1/3.3=10/33$.
Let $T\sim\text{Bin}(500,p_3)$.
Use normal approximation with continuity correction.
\end{scratch}

\textbf{(a)} $\mu_S=500\cdot\frac{5}{33}=\frac{2500}{33}\approx 75.76$ and
$\sigma_S^2=500\cdot\frac{5}{33}\cdot\frac{28}{33}\approx 64.28$, so $\sigma_S\approx 8.02$.
With continuity correction,
\[
P(S<65)=P(S\le 64)\approx \Phi\!\left(\frac{64.5-75.76}{8.02}\right)
=\Phi(-1.40)\approx 0.080.
\]

\textbf{(b)} $\mu_T=500\cdot\frac{10}{33}=\frac{5000}{33}\approx 151.52$ and
$\sigma_T^2=500\cdot\frac{10}{33}\cdot\frac{23}{33}\approx 105.6$, so $\sigma_T\approx 10.28$.
With continuity correction,
\[
P(T>160)=P(T\ge 161)\approx 1-\Phi\!\left(\frac{160.5-151.52}{10.28}\right)
=1-\Phi(0.87)\approx 0.191.
\]

\finalanswer{P(S<65)\approx 0.080,\ \ P(T>160)\approx 0.191}
\end{solution}

\newpage

\problem{4.26}
100 randomly chosen individuals were interviewed to estimate the unknown fraction $p \in (0, 1)$ of the population that prefers whole milk to skim milk. The resulting estimate is $\hat{p}$. With what level of confidence can we state that the true $p$ lies in the interval $(\hat{p}-0.1,\hat{p}+0.1)$?

\begin{solution}
\begin{scratch}
Let $S\sim\text{Bin}(n,p)$ with $n=100$, and $\hat p=S/n$.
We want the confidence level $P(|\hat p-p|<0.1)$.
Normal approximation gives
$P(|\hat p-p|<\varepsilon)\approx 2\Phi\!\left(\dfrac{\varepsilon\sqrt{n}}{\sqrt{p(1-p)}}\right)-1$.
Since $p(1-p)\le 1/4$, a lower bound is $2\Phi(2\varepsilon\sqrt{n})-1$.
\end{scratch}

With $\varepsilon=0.1$ and $n=100$,
\[
P(|\hat p-p|<0.1)\approx 2\Phi(0.2\sqrt{100})-1=2\Phi(2)-1\approx 0.9545.
\]
So the confidence level is about $95.45\%$.

\finalanswer{\text{confidence}\approx 0.9545}
\end{solution}


\end{document}
